\subsection{Sprint 2}

Iniciei o trabalho realizando diversas alterações no diagrama de arquitetura, conforme as solicitações do arquiteto. Durante esse processo, também precisei lidar com algumas particularidades envolvendo nossa equipe, como o fato de alguns colegas trabalharem como “heróis”, enquanto outros tentavam se esquivar de suas atribuições. Em paralelo, ofereci suporte aos AGES I e II na configuração do ambiente Docker, algo que gerou problemas até o final do projeto.

Outro ponto importante foi a preparação do ambiente na AWS, incluindo a criação da nossa instância EC2 e os testes de conexão via SSH. Esse processo exigiu uma série de depurações, já que inicialmente houve falhas ao tentar conectar-se ao servidor. Continuei investigando o problema até identificar a causa: embora o \textit{security group} da máquina estivesse aparentemente com tudo liberado, a conexão só funcionou após liberar explicitamente a porta 22 para SSH. Com o acesso finalmente resolvido, avancei para os testes iniciais de \textit{deploy} das aplicações \textit{frontend} e \textit{backend}, iniciando a validação do ambiente e preparando o caminho para as próximas etapas do desenvolvimento.
\\

\fbox{%
    \parbox{0.95\linewidth}{%
        \setlength{\parindent}{15pt}
            \itshape Infelizmente, nosso trabalho na infraestrutura e na arquitetura só pôde começar a essa altura, por fatores que posso atribuir tanto à nossa equipe quanto à equipe da AGES. Até esse ponto do semestre, não havíamos recebido orientações sobre onde encontrar instruções relativas ao orçamento da AWS e à criação do repositório da \textit{Wiki}, nem havíamos tido o acompanhamento do arquiteto com nosso time.

            Isso causou confusão entre nós, alunos de AGES III, que, depois desse tempo, procuramos o arquiteto da AGES para sanar nossas dúvidas. Foi somente nesse momento que descobrimos que, essencialmente, já deveríamos saber o que fazer e onde obter as informações necessárias.

            Senti muita falta de direcionamento. Por mais que exista o argumento de que nós, alunos, temos de ser proativos, isso pouco adianta quando nos faltam instruções mínimas. Se tivesse ocorrido uma conversa com nosso time no começo do semestre, mesmo que fosse apenas uma mensagem por e-mail, definitivamente teríamos adiantado muita coisa nessa \textit{sprint} e, consequentemente, no restante do semestre.
    }%
}
